\subsection{Teilmengen}
\label{subsec:fundamentals_subsets}

\begin{DefinitionBox}{Teilmenge}{fundamentals_subset}
  \begin{minipage}[t]{0.62\linewidth}
    \vspace{0pt}
    Eine Menge $A$ heißt \emph{Teilmenge} \EnglishTerm{subset} einer Menge
    $B$, wenn jedes Element von $A$ auch Element von $B$ ist. Man schreibt
    \[
      A\subseteq B
      \quad\Longleftrightarrow\quad
      \forall x:\; x\in A\Rightarrow x\in B.
    \]
  \end{minipage}\hfill
  \begin{minipage}[t]{0.32\linewidth}
    \vspace{0pt}
    \centering
    \resizebox{\linewidth}{!}{%
      \begin{tikzpicture}
  % Die äußere Fläche stellt B dar, die vollständig darin liegende Fläche A.
  \fill[THMBlue!16]
    (0,0) ellipse (3.8cm and 1.75cm);
  \draw[THMBlue!75!black,line width=1pt]
    (0,0) ellipse (3.8cm and 1.75cm);

  \fill[THMBlue!5]
    (0.85,-0.05) ellipse (1.75cm and 0.95cm);
  \draw[THMBlue!75!black,line width=1pt]
    (0.85,-0.05) ellipse (1.75cm and 0.95cm);

  \node[font=\Large,text=THMBlue!70!black] at (-1.85,0.1) {$B$};
  \node[font=\Large,text=THMBlue!70!black] at (0.85,-0.05) {$A$};
\end{tikzpicture}
%
    }
  \end{minipage}
\end{DefinitionBox}

\begin{RemarkBox}[title=Teilmenge und echte Teilmenge]
  \renewcommand{\arraystretch}{1.35}
  \begin{tabularx}{\linewidth}{@{}>{\bfseries}l c X@{}}
    \toprule
    Begriff & Schreibweise & Bedeutung \\
    \midrule
    Teilmenge
      & $A\subseteq B$
      & Jedes Element von $A$ ist auch Element von $B$; $A=B$ ist erlaubt. \\
    Echte Teilmenge
      & $A\subsetneq B$
      & $A\subseteq B$ und zusätzlich $A\neq B$. \\
    \bottomrule
  \end{tabularx}
\end{RemarkBox}


\begin{ExampleBox}{Teilmengen}{fundamentals_subset_example}
  Für $A=\{1,2\}$ und $B=\{1,2,3\}$ gilt
  \[
    A\ \Cloze{1.8cm}{\subseteq B},
    \qquad
    B\ \Cloze{1.8cm}{\not\subseteq A}.
  \]
\end{ExampleBox}
